% GENERATED from content/data/enrichment.json by mxbuild.enrichment; do not edit.
\mxunit{enrichment-invariants-exercises}{Parité et invariants}
\mxsubtitle{Seconde · raisonnement \textperiodcentered{} Exercices progressifs \textperiodcentered{} 01}

\begin{exercise}
On part de 7 et on ajoute ou retranche 4 à chaque coup. Peut-on atteindre 20 ? Et 23 ?
\end{exercise}

\begin{exercise}
Un produit de cinq nombres égaux à 1 ou $-1$ vaut $-1$. Montrer qu’il est impossible de rendre tous les nombres positifs en changeant deux signes à chaque coup.
\end{exercise}

\begin{exercise}
On remplace deux entiers a et b par $a + b$. Quel invariant conserve cette opération ?
\end{exercise}

\begin{exercise}
Pourquoi deux états ayant la même parité ne sont-ils pas nécessairement reliés par les coups autorisés ? Donner un exemple.
\end{exercise}

